\section{Future Work}\label{sec:future-work}

The most natural language extension is richer musical context.
Key signatures are a good example because they would let Motivo understand scale degrees, diatonic transposition, and
harmonic functions instead of treating all notes as independent pitch literals.
This should be introduced carefully.
It would require syntax, semantic rules, interactions with accidentals, and clear lowering behavior.
Because that support is not part of the current implementation, it belongs in future work rather than in the language
description.

Another useful extension is continuous musical data.
MIDI control changes, pitch bend, and automation curves would make Motivo more expressive for modern production
workflows.
Supporting them would require changes to the AST, type system, intermediate representation, MIDI backend, piano-roll
visualizer, and playback scheduler.
The existing separation between IR and backend makes such an extension plausible, but it is not a small addition.

The compiler could also expose more evaluation and debugging data.
For example, Motivo Studio could show the number of generated events, the longest track duration, compile time, MIDI
file size, and event-density warnings after each compile.
This would turn the current thesis measurements into user-facing feedback and would help users understand when a source
program is becoming too large.

Finally, the application could support project persistence and export workflows.
Saving Motivo source files, exporting MIDI with metadata, comparing generated versions, and integrating with cloud
storage would move Motivo Studio closer to a practical composition tool.
Those features would not change the compiler's core idea, but they would improve the surrounding product.

A standard library is another natural direction. Reusable scales, chord progressions, arpeggiators, Euclidean
rhythms, and common drum patterns would make the language more practical without changing its core compiler
model. Such a library should be designed after the core type system is stable, because library functions
would depend on how Motivo represents harmonic context and pattern parameters.

Alternative backends also remain possible future work. The current backend targets Standard MIDI Files, but
the same intermediate representation could be used as the basis for JSON event export, MusicXML-oriented
notation export, or integration with synthesis systems. Those targets would require separate mapping rules
and should not be treated as implemented behavior in the present thesis.
